% Replace only the original llm.c project section. Uses existing resume macros.
\projectHeading
{基于 llm.c 的 Mini-vLLM 推理引擎}
{C/C++ / OpenMP}

\begin{itemize}
    \item \textbf{增量解码链路：}
    基于 llm.c 扩展 GPT-2 单 Token 前向，复用模型权重及基础算子，
    按绝对位置更新位置编码并缓存各层历史 K/V；
    设计独立推理 Workspace，并在 Greedy Decoding 中跳过 Softmax。

    \item \textbf{分页缓存管理：}
    实现固定大小 CPU KV Cache Block Pool、逻辑 Block Table 与 Block Manager，
    支持缓存块按需分配、请求完成释放、跨请求复用及异常状态检测。

    \item \textbf{请求调度与 Paged Attention：}
    参考 vLLM 控制面设计 Sequence 状态及 Token Budget 驱动的 Scheduler，
    通过 GPT2ModelRunner 将调度结果转换为异长动态微批次，
    打通 Chunked Prefill/Decode、Greedy Sampling、状态提交及 Block 回收闭环；
    实现页表寻址的 K/V 写入、历史注意力计算及 Softmax 加权聚合，
    通过 OpenMP 并行处理 batch 与 attention head；
    在多层、多头、跨页及非连续页映射场景下与独立稠密参考对齐；
    构造 Decode/Prefill 混合及请求动态加入/退出场景，
    验证跨页扩容、释放复用，并与完整前缀 GPT-2 的全词表 logits 及生成 Token 对齐。
\end{itemize}
